\scanpage{11}%
\deleted{Noter d'autre part qu'on a équivalence entre les conditions suivantes:}

\deleted{a) L'homomorphisme $\varphi$ est surjectif (resp. injectif) pour tout $d$ et tout point $s\in U_d$ (définissant un $Y$ correspondant).}

\deleted{b) Même condition, mais avec \uline{un} $d$ fixé multiple de $p$, et tel que \deleted{\unsure{c)}} l'on ait $\mathrm{Ker}\,v=0$ (condition sans doute automatiquement \unsure{vérifié|vérifiée} (8.5.7)).}

\deleted{c) $(\mathrm{Tors}^{n-1}(X)\otimes_W k)_{ss}=0$ (resp. $(\mathrm{Tors}^{\handwritten{n+1}}(X)\otimes_W k)_{ss}=0$ ).}

\inserted{8.5.\deleted{11}\handwritten{\unsure{10|11}}.} Il serait intéressant de donner des exemples où \deleted{ces conditions ne sont \unsure{pas} satisfaites. Noter que la condition non respée est toujours satisfaite si $n=2$ ($X$ une surface)? Dans le cas général,} \inserted{\handwritten{$\varphi$ n'est pas bijectif}} on a l'inégalité
\[
\dim_k H^{n-1}(X,\underline O)_{ss}=\dim_{\underline{\mathrm F}_p}H^{n-1}(\overline X,\underline Z/p\underline Z)\handwritten{\geq}\dim\bigl((H^{\handwritten{n-1}}(X)/\mathrm{Tors}^{n-1}(X))\otimes k\bigr)_{ss}
\tag*{\inserted{(8.5.1\deleted{1}0.1)}}
\]
(grâce à 8.1.2 et 8.4.3 ; en fait cela marche en toute dimension $i$, pas seulement $i=n-1$), et dire que $\varphi$ est bijectif signifie que l'inégalité précédente est une égalité. L'exemple de Serre au symposium de Mexico fournit un $X$ (quotient d'une intersection complète par un groupe $G=\underline Z/p\underline Z$ opérant sans points fixes) pour lequel on a $H^1(X)=0$ , mais $H^1(X,\underline Z/pZ)\unsure{\simeq|\doteqdot|=|\ne}\underline Z/\unsure{p\underline Z}$. Quitte à couper avec des hpyperplans, on trouve le même exemple avec $\dim X=2$ . Pour cet exemple, \deleted{(8.5.11.1) n'est pas une égalité, et par suite,} d'après \deleted{ce qui précède,} \inserted{8.1.2, on aura} $(\handwritten{{}_p}\mathrm{Tors}^{\handwritten{2}}(X)\deleted{\unsure{\otimes_W k}})_{ss}\simeq\overset{\handwritten{H^1(X,\underline O_{\unsure{X|Y}})_{ss}}}{\overset{\handwritten{\shortparallel}}{H^1(X,\underline Z/p\underline Z)\otimes k}}=\mathrm{Ker}\,\handwritten{\varphi}$ (quel que soit $d$ et la section $Y$ de $X$ de degré $d$, permettant de définir $\varphi$ ) , qui est non nul, donc $\varphi$ n'est pas injectif. \deleted{Si maintenant $Z$ est un schéma lisse projectif sur $k$ de dimension $n-2$ ($n$ donné $\geq\deleted{\unsure{5|3}}\handwritten{4}$), à cohomologie cristalline sans torsion \inserted{et $(H^{n-4}(Z)\otimes k)_{ss}\ne0$} (par exemple $\mathrm P^{n-2}_k$ \inserted{si $n$ pair ...}), alors, posant $X'=X\times Z$ , on voit (sous réserve de la \uline{formule de Kunneth} en cohomologie cristalline) que $(\mathrm{Tors}^{n-1}(X')\otimes k)_{ss}$\deleted{\unsure{ntient comme facteur direct}} est non nul, ce qui fournit des exemples où $\varphi$ n'est pas surjectif, pour tout $n\geq4$. Il y a sans doute aussi des exemples avec $n=3$, en cherchant bien ....}
% archive sheet number: 10

